\scanpage{22}
\uline{Corollaire} 2 Sous les conditions du lemme\inserted{On suppose que
$\mathbb M_n(C_0)=\mathbb M_n(C_0)^*+\operatorname{Im}(\mathbb M_n(B)\to\mathbb M_n(C_0))$
[ce qui est vérifié sous les conditions signalées au N.B du Lemme 1]}, soit
$n\geq1$ un entier, $\mathbb M_n$ le foncteur ``matrices
carrées de degré $n$ à \unsure{cff. ds|coeff. ds|coeff. dans}'', $\mathrm{Gl}_n=\mathbb M_n^*$, alors on
a une bijection respectant l'unité
\[
\operatorname{Ker}\big(H^1(A'/A,\mathrm{Gl}_n)\to H^1(A'_0/A_0,\mathrm{Gl}_n)\big)\to
\operatorname{Ker}\big(H^1(A'/A,\mathbb M_n)\to H^1(A'_0/A_0,\mathbb M_n)\big)
\]
Par suite, si \deleted{$H^1(A'/A,G_a)=0$, alors}
\inserted{$H^1(A'_0/A_0,\mathrm{Gl}_n)=\{e\}$, $\underbrace{H^1(A'_0/A_0,\mathbb M_n)=\{e\}}_{\substack{\Updownarrow\\H^1(A'_0/A_0)=0}}$,}
on a une bijection
\[
H^1(A'/A,\mathrm{Gl}_n)\xrightarrow{\sim}H^1(A'/A,\mathbb M_n)=\deleted{\unsure{H^1(}}\mathbb M_n(H^1(A'/A))
\]
donc pour que $H^1(A'/A,\mathrm{Gl}_n)=\{e\}$, il f. et s. que
$H^1(A'/A,G_a)=\{e\}$.

Démonstration. Les deux mbs à comparer
sont égaux respectivement à
\[
H^1(I\deleted{\unsure{K}}C(A'/A,\mathbb M_n))/\mathrm{Gl}(n,C_0)
\]
et à
\[
H^1(I\deleted{K}C(A'/A,\mathbb M_n))/\mathbb M_n(C_0)
\]
or les opérations de $\mathrm{Gl}(n,C_0)$ sur $H^1(I\deleted{K}C(A'/A,\mathbb M_n))$
sont induites par celles de $\mathbb M_n(C_0)$. De plus
un élément de $\mathbb M_n(C_0)$ provenant de $\mathbb M_n(B)$ opère
trivialement. Donc il suffit encore de \deleted{\unsure{vérifier que}} \inserted{utiliser l'hypothèse}
% archive sheet number: 21
